% Vietnamese translation for OI Public Library
% Source-File: IOI中国国家候选队论文/2006/周以苏/周以苏.ppt
\documentclass[11pt]{article}
\usepackage{oipl-vn}

\title{Phản dịch hợp ngữ và tối ưu hằng số thời gian\\{\large Ghi chú theo slide}}
\author{Zhou Yisu}
\date{2006}

\begin{document}
\maketitle

\origtitle{Disassembly for time-constant optimization}
\sourcepath{IOI China candidate team papers / 2006 / Zhou Yisu / Zhou Yisu.ppt}

\section{Phạm vi}

Bộ slide đi kèm bài viết về tối ưu hằng số thời gian bằng cách quan sát mã
hợp ngữ do trình biên dịch sinh ra. Ghi chú này không thay thế phần bài viết
đầy đủ; nó gom lại các điểm nên nhớ khi đọc slide.

\section{Vì sao xét hằng số}

Trong thi lập trình, hai chương trình có cùng độ phức tạp tiệm cận vẫn có thể
khác nhau rất xa về thời gian chạy. Khi giới hạn dữ liệu sát với khả năng của
thuật toán, chi phí của phép chia, truy cập bộ nhớ, gọi hàm, truyền tham số,
hoặc kiểm tra nhánh đều có thể quyết định kết quả.

\section{Cách nhìn qua hợp ngữ}

Phản dịch hợp ngữ giúp người viết chương trình thấy một câu lệnh C++ được
biến thành bao nhiêu thao tác máy. Một phép gán mảng đơn giản thường khác một
phép truy cập qua con trỏ nhiều tầng; một vòng lặp có điều kiện phức tạp có
thể sinh thêm nhiều nhánh; phép chia nguyên thường đắt hơn cộng, trừ và dịch
bit.

\section{Cải tiến thực dụng}

Các hướng tối ưu trong bài viết gồm:
\begin{itemize}
  \item giảm phép toán đắt trong vòng lặp nóng;
  \item đưa biểu thức không đổi ra ngoài vòng lặp;
  \item chọn kiểu dữ liệu và cách truyền tham số phù hợp;
  \item thay đổi bố cục dữ liệu để truy cập bộ nhớ tuần tự hơn;
  \item kiểm tra mã sinh ra sau khi tối ưu để tránh sửa nhầm vào chỗ không
  ảnh hưởng đáng kể.
\end{itemize}

\section{Giới hạn}

Tối ưu hằng số chỉ có ý nghĩa sau khi thuật toán lớn đã đúng hướng. Nếu độ
phức tạp tiệm cận sai, việc viết sát máy chỉ che được một phần rất nhỏ. Ngược
lại, khi mô hình đã đúng và nút thắt nằm trong cài đặt, đọc mã hợp ngữ giúp
ta đưa ra quyết định dựa trên bằng chứng thay vì cảm giác.

\end{document}
